\begin{resumo}[Abstract]
The study of Theory of Computation is essential for students in computer-related fields, providing the theoretical foundation required to design modern compilers. However, concepts associated with Context-Free Languages (CFLs) and their recognizers, Non-Deterministic Pushdown Automata (NPDAs), frequently pose significant learning barriers due to their high level of abstraction and the inherent complexity of tracking multiple parallel computations. This work presents the development of a web-based educational tool designed for the graphical modeling and real-time simulation of NPDAs. The methodology consisted of developing a client-side (\textit{frontend}) application using vanilla JavaScript and the Cytoscape library for graph rendering. The system enables users to design automata, define distinct acceptance criteria, and dynamically analyze the computation branches generated by non-determinism through structures called \textit{snapshots}. As a result, the platform allows exporting and importing models in JSON format, as well as managing exercises in a completely synchronous and free environment. Although single-thread processing introduces performance constraints when dealing with massive data volumes, the tool successfully fulfills its pedagogical purpose of mitigating traditional static visualization obstacles. Future works include conducting empirical evaluations with students to validate its pedagogical effectiveness, expanding the scope to cover other state machines, and architecting a \textit{backend} for performance optimization.

\textbf{Keywords}: Theory of Computation. Non-Deterministic Pushdown Automata. Educational Tool. Algorithm Visualization. Web Platform.
\end{resumo}